\chapter{Literature Review}
\label{ch:literature}

\section{Existing Approaches}

\subsection{Conventional control}
Fixed-time plans set the cycle length and green splits from historical flows, and actuated
systems extend or end a green phase from detector calls with hand-written rules.
Max-Pressure \citep{varaiya2013max} gives green, at each intersection, to the phase with the
largest difference between upstream and downstream queues. It needs only local measurements,
and it keeps the queues of the whole network stable whenever demand lies inside the capacity
region, in a model with unlimited lane storage and known turning ratios. Rule controllers
are transparent and need no training, but the rule never changes. The stability result has
two limits that matter for this project: it is about mean queue lengths, not the wait of a
single vehicle, and its model has neither finite lanes nor time lost at phase changes.

\subsection{Single-agent reinforcement learning}
Deep Q-networks (DQN) \citep{mnih2015human} learn the value of each action from a state
vector and were soon applied to isolated intersections. A single agent either controls one
intersection or treats every intersection on its own with shared weights, which is the
single-agent baseline used here. These agents output a phase and nothing about how far that
choice can be trusted.

\subsection{Multi-agent reinforcement learning}
The simplest multi-agent form gives every intersection its own independent learner (IDQN),
although each agent then faces an environment that changes as the others learn. PressLight
\citep{wei2019presslight} uses the negative pressure as the reward, which ties the learned
objective to Max-Pressure. CoLight \citep{wei2019colight} adds graph attention over the
neighbouring intersections and released the Jinan and Hangzhou flows used in this project.
The gains reported for newer methods are fragile: in a large benchmark a well-tuned
independent DQN matched most of them \citep{ault2021reinforcement}. Two habits limit what
published numbers show. Studies differ in decision interval, transition time and flow files,
so their travel times cannot be compared directly, and evaluation is usually done with perfect
detectors on the demand the agent was trained for, often without reporting the spread
between training runs. Most recent work uses the CityFlow simulator
\citep{zhang2019cityflow}, as this project does.

\subsection{Uncertainty and safe reinforcement learning}
Epistemic uncertainty comes from limited training data, unlike the randomness of the traffic
itself, and an ensemble of independently initialised networks is a simple way to estimate it.
Bootstrapped DQN \citep{osband2016deep} trains several Q-heads on one shared network and
gives each head a random subset of the experience. In signal control, ensembles like this
have been used to explore during training \citep{tan2020multiagent}, not to judge a decision
once the controller is running. In automated driving, by contrast, Hoel et al.\
\citep{hoel2020tactical} let an ensemble's disagreement switch a vehicle to a safe action at
execution time. The concurrent Com-MP \citep{zhu2026commp} combines RL with Max-Pressure by
keeping Max-Pressure as the phase selector and letting RL adapt it, which keeps a stability
claim by construction but confines learning to the rule.

\subsection{Explainability}
Attention weights are often shown as explanations of a controller's decisions
\citep{wei2019colight}. A weight is not the same as an influence on the output, though, so
whether attention explains a given model has to be tested on that model.

\section{Research Gaps}

Three gaps stand out against the requirements of the problem statement.

\textbf{No reliability signal when the controller acts.} RL signal controllers return a phase
and no sign of whether the decision can be trusted. Where ensembles of Q-heads exist, their
disagreement serves exploration during training and is discarded afterwards
\citep{tan2020multiagent}.

\textbf{Performance and guarantees live in different controllers.} Learned controllers are
best on average and have no guarantee; Max-Pressure has a stability result and does not
learn. Falling back to a safe controller when an ensemble disagrees is established for a
single vehicle \citep{hoel2020tactical}, and the one recent hybrid of RL and Max-Pressure
leaves the rule in charge \citep{zhu2026commp}. The search behind this review
found no multi-agent signal controller that lets a learned policy lead and hands individual
decisions back to Max-Pressure when its uncertainty is high.

\textbf{Robustness and explanations are rarely tested.} Most studies evaluate with perfect
detectors on the training demand, and attention maps are presented without checking that they
reflect what the network uses.

\ugmarl{} answers the first two gaps with a gate that reads ensemble disagreement while the
controller runs, and the third with tests under detector noise, dropout, demand surge and an
unseen city, plus a faithfulness check of the attention weights. Two published methods are
closest. Com-MP also joins learning with Max-Pressure, but Max-Pressure always selects the
phase there, whereas in \ugmarl{} the learned policy selects it and Max-Pressure takes over
only when the uncertainty is high. Hoel et al.\ use the same chain of ensemble, threshold
and fallback for one vehicle; here it runs at every intersection of a multi-agent network,
with a fallback that has a stability result. The literature search was not exhaustive, so
this claim holds only as far as that search reached. Attention-based coordination \citep{wei2019colight}, the
pressure reward \citep{wei2019presslight}, bootstrapped Q-heads \citep{osband2016deep} and the
stability of Max-Pressure \citep{varaiya2013max} come from the cited work and are not claimed
as new.
